\documentclass[11pt]{article}
\usepackage[a4paper,margin=1in]{geometry}
\usepackage{amsmath,amssymb,mathtools,bm}
\usepackage{enumitem,booktabs,array}
\usepackage{tcolorbox}
\usepackage{hyperref}
\usepackage[protrusion=true,expansion=false]{microtype}
\hypersetup{colorlinks=true,linkcolor=blue,urlcolor=blue,citecolor=blue}
\setlist[itemize]{topsep=2pt,itemsep=2pt}
\setlist[enumerate]{topsep=2pt,itemsep=2pt}
\newtcolorbox{keybox}{colback=blue!3!white,colframe=blue!40!black,boxrule=0.4pt,arc=1mm}
\newtcolorbox{alertbox}{colback=red!3!white,colframe=red!40!black,boxrule=0.4pt,arc=1mm}
\newtcolorbox{answerbox}{colback=green!3!white,colframe=green!40!black,boxrule=0.4pt,arc=1mm}
\newtcolorbox{drillbox}{colback=yellow!5!white,colframe=orange!60!black,boxrule=0.4pt,arc=1mm}
\newcommand{\EV}{\mathbb{E}}
\newcommand{\Var}{\mathrm{Var}}
\newcommand{\Cov}{\mathrm{Cov}}
\newcommand{\blank}[1]{\underline{\hspace{#1}}}

\title{E03-02: Ahern \& Dittmar (2012, QJE) \\
\large ``The changing of the boards: The impact on firm value of mandated female board representation'' --- Active Recall Workbook}
\author{Empirical Corporate Finance II --- Exam Prep}
\date{\today}
\begin{document}\maketitle

\begin{keybox}
\textbf{Paper in one sentence.} Norway's 2003--2008 gender quota (40\% female board representation for public limited firms) caused a sharp decline in Tobin's Q --- roughly $-3.5$\% at announcement and deeper over implementation --- driven by mandatory appointment of younger, less-experienced female directors, documenting a real cost of rigid diversity mandates.

\textbf{Placement.} Canonical early study of the Norwegian quota; launched the literature on board diversity and forced representation. Paired with Eckbo-Nygaard-Thorburn (E03-03) who revisit with cleaner identification and find the effect is much smaller.
\end{keybox}

\section*{Round 1: Complete Walk-Through}

\subsection*{1.1 Setting}
Norway's ASA-law (2003): requires 40\% female representation on all public limited (ASA) firm boards by 2008. Non-compliance risks forced dissolution. Unique exogenous shock to board composition.

\subsection*{1.2 Sample}
\begin{itemize}
\item $\sim$250 Norwegian ASA firms.
\item Private limited (AS) firms not bound by quota; serve as control.
\item Event windows: 2003 announcement, 2006 legal binding, 2008 full compliance.
\end{itemize}

\subsection*{1.3 Empirical design}
Event study around announcement plus DiD comparing ASA vs.\ AS firms:
\[
\Delta Q_{it} = \alpha + \beta\cdot\text{PostQuota}_{t}\cdot\text{ASA}_{i}+X_{it}'\gamma+\varepsilon_{it}.
\]
Firm FE and industry-year controls.

\subsection*{1.4 Headline results}
\begin{itemize}
\item Announcement CAR for ASA firms: $\approx -3.5$\%.
\item Tobin's Q declines $\sim$12--15\% over implementation.
\item Effect strongest for firms with \emph{least} female-directorship slack (i.e.\ had to do most adjustment).
\item Firms with more diverse pre-quota boards largely unaffected.
\end{itemize}

\subsection*{1.5 Mechanism}
\begin{itemize}
\item Quota forced rapid appointment of women who were younger, less experienced, less connected.
\item Average tenure and board experience of new female directors significantly below replaced male directors.
\item Suggests the \emph{cost of forced search} under tight deadlines, not of female directors per se.
\end{itemize}

\subsection*{1.6 Interpretation}
\begin{itemize}
\item Rigid quotas impose adjustment costs when the pool of qualified candidates is shallow.
\item Value effect is an implementation phenomenon, not a critique of diversity.
\item Policy implication: gradual phase-ins and development of talent pipelines matter.
\end{itemize}

\subsection*{1.7 Caveats}
\begin{alertbox}
\begin{itemize}
\item AS firms differ systematically from ASA; control inference is arguable.
\item Short post-implementation window.
\item Later work (Eckbo-Nygaard-Thorburn 2022) substantially revises the magnitudes --- see E03-03.
\end{itemize}
\end{alertbox}

\section*{Round 2: Level 1 Fill-in}
\begin{enumerate}
\item Country / quota: \blank{1.5cm}, \blank{1cm}\% female representation.
\item Binding year: \blank{1cm}; full compliance: \blank{1cm}.
\item Announcement CAR: $\sim$\blank{1cm}\%.
\item Mechanism: new female directors were \blank{1.5cm}/less experienced.
\item Paper that revisits: Eckbo, Nygaard \& \blank{2cm} (2022).
\end{enumerate}

\section*{Round 3: Level 2 Prompts}
\begin{enumerate}
\item Describe the Norwegian ASA quota and why it gives a clean setting.
\item Write the DiD spec and discuss identification.
\item Report the headline results and the proposed mechanism.
\item How does AD interpret the value decline (implementation vs.\ directors themselves)?
\item Discuss policy implications and relate to Eckbo-Nygaard-Thorburn (2022).
\end{enumerate}

\section*{Round 4: Minimal Scaffolding}
From a blank page: (i) setting, (ii) DiD, (iii) value effects, (iv) mechanism, (v) interpretation.

\section*{Round 5: Blank Page}
Essay: board quotas and firm value --- the AD (2012) evidence.

\vspace{6pt}
\begin{drillbox}
\textbf{Speed Drill.} (1) Quota year and \%. (2) Announcement CAR. (3) Q decline. (4) Mechanism. (5) Follow-up paper.
\end{drillbox}

\section*{Quick Reference \& Answer Key}
\begin{answerbox}
\begin{itemize}
\item \textbf{Policy.} Norway 40\% quota, 2003--08.
\item \textbf{CAR.} $\approx -3.5$\%.
\item \textbf{Q.} Down $\sim$12--15\%.
\item \textbf{Mechanism.} Forced appointment of younger, less-experienced directors.
\end{itemize}
\end{answerbox}

\begin{keybox}
\textbf{30-second exam pitch.} Ahern \& Dittmar (2012) study Norway's 2003 gender quota requiring public limited (ASA) firm boards to be at least 40\% female by 2008. They document a sharp Tobin's Q decline at treated firms: announcement CAR $\approx -3.5$\% and a $\sim$12--15\% Q decline through implementation, concentrated at firms with the least pre-quota diversity (thus the largest adjustment). The mechanism is that newly-appointed female directors were disproportionately younger, less-experienced, and less-connected than the male directors they replaced --- a forced-search cost, not an indictment of diversity per se. The paper launches the board-diversity and forced-representation literature, raises policy concerns about rigid timelines and shallow talent pipelines, and should be read alongside Eckbo-Nygaard-Thorburn (2022), which substantially revises the magnitude downward with cleaner identification.
\end{keybox}

\end{document}
